\section{Modul 03. Runtunan}

\subsection{Tujuan}

Setelah menyelesaikan modul ini, mahasiswa mampu:
\begin{enumerate}[nosep]
    \item menjelaskan tugas instruksi \texttt{MOV}, \texttt{XCHG},
          \texttt{ADD}, \texttt{SUB}, \texttt{INC}, dan \texttt{DEC};
    \item menulis runtunan instruksi yang dijalankan dari atas ke
          bawah tanpa percabangan dan tanpa pengulangan;
    \item menelusuri perubahan isi register setelah setiap instruksi;
    \item mengubah satu digit 0--9 menjadi kode ASCII dengan
          \texttt{ADD AL,30h} untuk dicetak;
    \item merakit, menjalankan, dan memeriksa register program dengan
          emu8086 atau DEBUG.
\end{enumerate}

\subsection{Dasar singkat}

\subsubsection*{Pemindahan data: MOV dan XCHG}

\texttt{MOV} menyalin nilai dari sumber ke tujuan. Nilai sumber tidak
berubah, sehingga instruksi ini disebut menyalin, bukan memindahkan.
Ukuran tujuan dan sumber harus sama: 8 bit dengan 8 bit, atau 16 bit
dengan 16 bit. Contohnya, \texttt{MOV DX,AX} menyalin isi \texttt{AX}
ke \texttt{DX}; \texttt{AX} tetap bernilai sama.

\texttt{XCHG} menukar isi dua operand. Setelah
\texttt{XCHG AX,BX}, isi \texttt{AX} berpindah ke \texttt{BX} dan isi
\texttt{BX} berpindah ke \texttt{AX}. Kedua instruksi tidak mengubah
flag hasil aritmatika. Perubahan register pada contoh program
ditunjukkan pada Tabel~\ref{tab:m03-pindah}.

\begin{table}[H]
    \centering
    \caption{Contoh perubahan register oleh \texttt{MOV} dan
    \texttt{XCHG}. Nilai awal: \texttt{AX=1234h}, \texttt{BX=5678h},
    \texttt{CX=9ABCh}.}
    \label{tab:m03-pindah}
    \small
    \begin{tabular}{@{}lll@{}}
        \toprule
        \textbf{Instruksi} & \textbf{Sesudah} & \textbf{Keterangan} \\
        \midrule
        \texttt{MOV DX,AX} & \texttt{DX=1234h} & menyalin, \texttt{AX}
        tetap \\
        \texttt{XCHG AX,BX} & \texttt{AX=5678h}, \texttt{BX=1234h}
        & menukar \\
        \texttt{XCHG AX,CX} & \texttt{AX=9ABCh}, \texttt{CX=5678h}
        & menukar \\
        \bottomrule
    \end{tabular}
\end{table}

\subsubsection*{Aritmatika: ADD, SUB, INC, dan DEC}

\texttt{ADD} menambahkan sumber ke tujuan dan menyimpan hasilnya pada
tujuan. \texttt{SUB} mengurangkan sumber dari tujuan. Keduanya
memerlukan ukuran operand yang sama. \texttt{INC} menambah 1 pada
satu operand, sedangkan \texttt{DEC} mengurangi 1. Karena hanya
memakai satu operand, \texttt{INC} dan \texttt{DEC} lebih ringkas
daripada \texttt{ADD} atau \texttt{SUB} dengan nilai 1.

Runtunan dijalankan berurutan dari atas ke bawah. Setiap instruksi
bekerja pada nilai yang ditinggalkan instruksi sebelumnya. Contoh
perubahan \texttt{AL} pada program kedua ditunjukkan pada
Tabel~\ref{tab:m03-hitung}.

\begin{table}[H]
    \centering
    \caption{Perubahan \texttt{AL} pada runtunan aritmatika.}
    \label{tab:m03-hitung}
    \small
    \begin{tabular}{@{}cll@{}}
        \toprule
        \textbf{Langkah} & \textbf{Instruksi} & \textbf{Nilai \texttt{AL}} \\
        \midrule
        1 & \texttt{MOV AL,5} & 5 \\
        2 & \texttt{ADD AL,3} & 8 \\
        3 & \texttt{SUB AL,1} & 7 \\
        4 & \texttt{INC AL} & 8 \\
        5 & \texttt{DEC AL} & 7 \\
        6 & \texttt{ADD AL,30h} & 37h, karakter \texttt{'7'} \\
        \bottomrule
    \end{tabular}
\end{table}

\subsubsection*{Mencetak satu digit}

Kode ASCII untuk angka 0 sampai 9 adalah \texttt{30h} sampai
\texttt{39h}. Karena itu, satu digit 0--9 dapat diubah menjadi
karakter dengan menambahkan \texttt{30h}. Hasilnya diletakkan pada
\texttt{DL}, lalu dicetak dengan \texttt{AH=02h}. Cara ini hanya
berlaku untuk satu digit. Hasil 10 atau lebih tidak dapat dicetak
dengan penambahan \texttt{30h}; nilainya harus diperiksa pada
register dengan emu8086 atau DEBUG. Modul ini tidak membuat rutin
pencetak bilangan banyak digit.

Runtunan pada modul ini tidak memakai \texttt{CMP}, lompatan, atau
\texttt{LOOP}. Semua instruksi dijalankan berurutan sampai program
berakhir.

Rujukan singkat: Gunawan, Bagian 5 (MOV dan XCHG); S'to, Bab VIII
(ADD, SUB, INC, dan DEC); praktikum \texttt{bab1.tex} dan
\texttt{bab2.tex}; buku ajar, \texttt{transfer-data.asm}.

\subsection{Alat dan bahan}

\begin{itemize}[nosep]
    \item Komputer dengan emu8086 terpasang.
    \item DOSBox dengan TASM, TLINK, dan DEBUG.
    \item Berkas \texttt{code/m03-pindah-data.asm} dan
          \texttt{code/m03-hitung.asm}.
    \item Lembar jurnal praktikum.
\end{itemize}

\subsection{Contoh program}

Program pertama memindahkan data antar register dengan \texttt{MOV}
dan \texttt{XCHG}. Program tidak mencetak keluaran; hasilnya
diperiksa pada jendela register emu8086 atau dengan DEBUG.

\lstinputlisting[
    caption={Pemindahan data dengan \texttt{MOV} dan \texttt{XCHG}.},
    label={lst:m03-pindah}
]{code/m03-pindah-data.asm}

Program kedua menjalankan \texttt{ADD}, \texttt{SUB}, \texttt{INC},
dan \texttt{DEC}. Satu hasil yang bernilai 7 dicetak sebagai karakter,
sedangkan hasil penjumlahan \texttt{1234h+5678h} diperiksa pada
register karena lebih dari satu digit.

\lstinputlisting[
    caption={Runtunan aritmatika dan pencetakan satu digit.},
    label={lst:m03-hitung}
]{code/m03-hitung.asm}

\subsection{Langkah kerja}

\begin{enumerate}
    \item Siapkan folder kerja, lalu salin kedua berkas \texttt{.asm}
          ke dalamnya.
    \item Rakit dan tautkan program pertama. Jalankan kedua perintah
          berikut pada baris terpisah:
          \begin{quote}
          \texttt{tasm m03-pindah-data.asm}\\
          \texttt{tlink m03-pindah-data.obj}
          \end{quote}
          Jalankan program itu di emu8086 dengan \emph{single step}.
    \item Catat isi \texttt{AX}, \texttt{BX}, \texttt{CX}, dan
          \texttt{DX} sesudah setiap instruksi. Bandingkan hasilnya
          dengan Tabel~\ref{tab:m03-pindah}. Jelaskan perbedaan kerja
          \texttt{MOV} dan \texttt{XCHG}.
    \item Rakit dan tautkan program kedua dengan perintah yang sama,
          lalu jalankan. Catat karakter yang tercetak.
    \item Periksa \texttt{AX} sesudah \texttt{ADD AX,5678h} dengan
          emu8086 atau DEBUG. Catat nilainya dan jelaskan mengapa
          hasil itu tidak dicetak.
    \item Ubah nilai awal pada program kedua sehingga runtunan
          menghasilkan digit lain antara 0 dan 9. Rakit ulang,
          jalankan, dan pastikan karakter yang tercetak berubah.
\end{enumerate}

\subsection{Tugas}

\begin{enumerate}
    \item Jelaskan perbedaan \texttt{MOV AX,BX} dan
          \texttt{XCHG AX,BX}.
    \item Diketahui \texttt{AX=1111h} dan \texttt{BX=2222h}. Tentukan
          isi kedua register sesudah \texttt{XCHG AX,BX}.
    \item Jelaskan tugas \texttt{ADD}, \texttt{SUB}, \texttt{INC},
          dan \texttt{DEC}, serta mengapa \texttt{INC} lebih ringkas
          daripada \texttt{ADD AL,1}.
    \item Jelaskan mengapa \texttt{ADD AL,30h} hanya dapat dipakai
          untuk mencetak satu digit 0--9. Tentukan hasilnya bila
          \texttt{AL} bernilai 10.
    \item Ubah program kedua agar runtunannya menghasilkan
          $9-3+1=7$, lalu tentukan nilai \texttt{AL} sebelum dan
          sesudah \texttt{ADD AL,30h}.
\end{enumerate}

\subsection{Hasil yang dikumpulkan}

\begin{enumerate}
    \item Kedua berkas \texttt{.asm} beserta versi modifikasinya.
    \item Tabel isi register sesudah setiap instruksi pada program
          pertama.
    \item Tangkapan layar keluaran program kedua dan tampilan
          register sesudah \texttt{ADD AX,5678h}.
    \item Lembar jurnal berisi jawaban langkah kerja dan tugas.
\end{enumerate}
